\documentclass{article}
\usepackage[utf8]{inputenc}
\usepackage{amsmath}
\usepackage{listings}
\usepackage{xcolor}
\usepackage{geometry}
\geometry{a4paper, margin=1in}

\lstset{
    basicstyle=\ttfamily\small,
    frame=single,
    breaklines=true,
    backgroundcolor=\color{gray!5},
    linewidth=0.7\textwidth,
    xleftmargin=0.1\textwidth
}

\title{Labwork 2: Linear Regression Implementation}
\author{Nguyễn Duy Anh Tùng - 2540107}
\date{May 8, 2026}

\begin{document}

\maketitle

\section{Introduction}
This report details the implementation of a Linear Regression model using the Gradient Descent algorithm. The objective is to find the optimal parameters $w_1$ and $w_0$ for the model $y = w_1 \cdot x + w_0$.

\section{Implementation}
The implementation focuses on minimizing the Mean Squared Error (MSE) loss function:
\begin{equation}
J(w_0, w_1) = \frac{1}{2N} \sum_{i=1}^{N} (w_1 \cdot x_i + w_0 - y_i)^2 
\end{equation}

To update the parameters simultaneously, partial derivatives were calculated as follows:
\begin{itemize}
    \item $\frac{\partial L}{\partial w_0} = w_1 \cdot x_i + w_0 - y_i$ 
    \item $\frac{\partial L}{\partial w_1} = x_i \cdot (w_1 \cdot x_i + w_0 - y_i)$ 
\end{itemize}

The parameters are updated iteratively: $w = w - r \cdot \frac{\partial L}{\partial w}$.

\section{Experimental Results}
The model was trained on the provided "lr.csv" dataset. Below are the final iterations showing convergence with a learning rate $r = 0.0008$:

\begin{center}
\begin{lstlisting}
...
37200      | 47.9727    | 1.2625     | 9.4513    
37300      | 47.9731    | 1.2625     | 9.4512
Final model: y = 1.2622x + 47.9878
\end{lstlisting}
\end{center}

\section{Learning Rate Analysis}
The learning rate $r$ is a critical hyperparameter:
\begin{itemize}
    \item \textbf{Optimal rate ($r = 0.0008$):} The model converges smoothly around epoch 37,000. In contrast, lower learning rates significantly increase the required number of epochs; for instance, at $r = 0.0001$, convergence takes approximately 300,000 epochs.
    \item \textbf{High rate:} Increasing $r$ beyond a certain threshold (somewhere around 0.0008 to 0.0009) leads to divergence, where the loss becomes infinite.
    \item \textbf{Low rate:} A smaller $r$ requires significantly more iterations to reach the same loss value.
\end{itemize}

\section{Conclusion}
The Linear Regression model was successfully implemented and optimized using 2D Gradient Descent. The resulting parameters provide a line that best fits the sample dataset.

\end{document}